\chapter{相关技术}

MAGIC4FDG系统的设计融合了多个领域的技术。本章介绍系统依赖的五项核心技术基础：模糊测试与模糊驱动的基本原理、LibFuzzer引擎的覆盖率引导机制、大语言模型的代码生成能力、多智能体协作框架，以及基于clang AST的静态分析。

这些技术看似分属不同领域，但在本系统中它们各司其职又紧密配合：模糊测试定义了任务目标和质量标准，覆盖率引导提供迭代反馈信号，大语言模型负责代码合成和约束推理，多智能体框架实现职责分离和协作编排，静态分析则提供结构化的上下文知识。

\section{模糊测试与模糊驱动}

模糊测试的历史可以追溯到1990年。Miller等人向90个Unix工具喂入随机字符串，发现约24\%的程序崩溃了，这种用随机输入触发异常的方法就此得名"fuzzing"\cite{miller1990}。它的核心思想在于尽可能广泛地探索程序的输入空间来暴露缺陷，几乎不需要人工编写测试用例。三十多年过去，模糊测试已经从纯随机输入演化为覆盖率引导的智能化方法，是当前发现安全漏洞最有效的技术之一\cite{Manès2021}。

从运行机制看，模糊测试围绕一个不断迭代的循环展开。引擎拿到初始种子后，对其施加各种变异，如位翻转、字节插入删除、整数边界替换等，并批量产出测试用例\cite{zhang2016}。这些用例被逐一送入经过插桩的被测程序执行。插桩代码在运行时记录覆盖率信息，同时系统监控是否出现崩溃或内存越界。覆盖率引导模式下的核心判据是：某个输入是否触发了此前未见的代码边。触发了就留下继续变异，没有就丢弃。值得一提的是，AFL和LibFuzzer等主流引擎采用的都是边覆盖而非路径覆盖\cite{hsu2018}。也就是说，它们关心控制流图中基本块之间的转移边有没有被走到，并不试图区分完整的执行路径序列。


在模糊测试系统中，三大核心组件各有分工：输入生成器基于种子做变异，变异质量直接决定探索效率；执行器把输入送入目标程序；缺陷监控器通过AddressSanitizer、UBSan等工具捕获异常，同时采集覆盖率反馈\cite{clusterfuzz2025}。通常，输入生成器和缺陷监控器合在一起构成模糊引擎，而执行器的核心就是模糊驱动。

模糊驱动的角色是把模糊引擎产生的原始字节流转化为被测库API能接受的合法调用序列\cite{Manès2021}。它需要完成输入解析、环境初始化、API调用和资源清理这几件事。以LibFuzzer为例，驱动要实现一个\texttt{LLVMFuzzerTestOneInput}函数，接收字节数据和长度，在函数体内把字节转化为有效的API调用\cite{libfuzzer2024}。

驱动质量对测试效果的影响是决定性的。好的驱动能把随机输入映射到深层代码路径，触发复杂行为；差的驱动让大量输入在浅层校验处就被拒绝。驱动本身的正确性也很关键，错误的参数初始化或遗漏前置条件会导致虚假崩溃或漏掉真实漏洞。目前工业界的驱动主要靠人工编写，耗时且覆盖有限\cite{lyu2024}。这正是本文研究自动生成的出发点。


\section{LibFuzzer模糊引擎与覆盖率引导技术}

LibFuzzer是LLVM项目提供的进程内、覆盖率引导、进化式模糊测试引擎\cite{libfuzzer2024}。和AFL这类进程外模糊器不同，LibFuzzer把自身和被测库链接成同一个二进制，在单进程内反复调用目标函数并变异输入，省去了进程创建销毁的开销，吞吐量很高。如图\ref{fig:driver-code}所示，它要求用户实现一个固定签名的目标函数，引擎运行时反复调用这个函数，每次传入不同的变异数据。

\begin{figure}[htbp]
\centering
\begin{lstlisting}[language=C, style=style@base]
   int LLVMFuzzerTestOneInput(const uint8_t *Data, size_t Size) {
       // 1. 输入数据解析与参数构造
       if (Size < 4) return 0;
       int flags = *(int*)Data;
       const char *payload = (const char*)(Data + 4);
       size_t payload_len = Size - 4;
       // 2. 环境初始化
       context_t *ctx = target_init(flags);
       // 3. 目标API调用
       target_process(ctx, payload, payload_len);
       // 4. 资源清理
       target_cleanup(ctx);
       return 0;
   }
\end{lstlisting}
\caption{LibFuzzer模糊驱动代码结构示例}
\label{fig:driver-code}
\end{figure}

LibFuzzer的覆盖率引导依赖LLVM的SanitizerCoverage插桩框架\cite{sanitizercoverage2024}。编译阶段，编译器在每个基本块入口插入回调代码，运行时记录哪些位置被执行。引擎内部用一个8位边计数器位图来跟踪状态：每条控制流边占一个槽位，记录的不是"是否走过"而是"命中次数落在哪个区间"。一旦某个新输入让位图出现了之前未见的模式（触发了新边，或让某条边的命中次数跳到新区间），引擎就将其保留继续变异。这种设计比单纯的覆盖/未覆盖二值判断更精细，因为它能区分循环执行一次和执行多次的行为差异。

实际使用中，LibFuzzer支持fork模式（\texttt{-fork=N}）来隔离崩溃。主进程管理语料库，依次启动子进程跑模糊测试；子进程独立运行一段时间后退出，把新覆盖输入汇报给主进程。某个输入触发崩溃只影响当前子进程，主进程随即启动新的继续探索。这对长时间运行很重要，单个崩溃不会中断对其他路径的探索。但fork模式也带来一个问题：工作进程被信号终止时，profraw文件可能没写完整。本系统通过语料库重放解决这个问题，fork模式结束后，用收集到的语料库在非fork模式下重跑一次（\texttt{-runs=0}），确保覆盖率数据完整。


精确的覆盖率数据收集依赖LLVM的llvm-cov工具链\cite{llvmcov2024}。具体做法是：编译阶段给目标加上\texttt{-fprofile-instr-generate -fcoverage-mapping}两个标志，启用源码级插桩。程序运行结束后会落盘一份.profraw文件，记录原始的计数器数据。随后用\texttt{llvm-profdata merge}将可能存在的多份profraw合并为统一的profdata格式。最终通过\texttt{llvm-cov export}导出JSON报告，报告中包含每一行的执行次数以及每个分支条件的真/假覆盖状态。

这种粒度的数据对本系统而言不可或缺。粗粒度反馈（比如只告诉你哪些函数被调用过、哪些API组合被覆盖了）无法定位具体的薄弱点。而行级加分支级的信息能够直接告诉缺口诊断智能体：某个\texttt{if}分支从未走过false路径，或者某段错误处理代码完全没被触及。有了这样明确的分析目标，后续的迭代改进才有据可依。

\section{大语言模型与代码生成技术}

大语言模型（LLM）是基于Transformer架构\cite{vaswani2023}的深度神经网络，通过在海量文本上预训练获得能力。Transformer最关键的创新是自注意力机制：处理序列中任何一个位置时都能直接关注其他所有位置，有效捕捉长距离依赖。这对代码生成尤为重要，因为变量定义和使用、函数声明和调用之间经常相隔数十甚至上百行，传统RNN难以建模这种远距离关联。从GPT\cite{radford2018}、BERT\cite{devlin2018}到GPT-4\cite{openai2023}，参数量从数亿增长到数千亿，"预训练+指令微调+RLHF"的训练范式让模型既具备丰富知识储备，又能遵循复杂指令\cite{yang2023}。

代码生成是检验LLM能力的一个重要维度。编程语言语法严格、语义明确，对模型的精确推理能力要求很高。OpenAI的Codex在HumanEval上达到28.8\%的单次通过率\cite{chen2021a}，此后GPT-4等模型持续提升。本文采用的Anthropic Claude系列\cite{anthropic2024claude}在代码生成和推理上表现优异，Claude Opus能在单次生成中产出数百行完整程序，多轮交互中对复杂约束的理解保持一致，这对本系统的多轮迭代生成至关重要。


在软件测试领域，LLM的应用越来越广。TitanFuzz\cite{deng2023b}证明了LLM可以零样本生成深度学习库的测试输入，Fuzz4All\cite{xia2024}把这个能力推广到任意语言。在驱动生成方面，PromptFuzz\cite{lyu2024}用迭代提示引导LLM持续生成新驱动，CKGFuzzer\cite{xu2024}用知识图谱增强模型对被测库的理解。这些工作表明LLM确实具备理解API语义、推断调用约束并生成规范测试代码的能力。

不过LLM用于代码生成并非没有代价。最突出的问题是"幻觉"：模型有时会自信地写出根本不存在的API名称，或者给函数传入类型不匹配的参数。另一个现实约束来自上下文窗口：当需要同时塞入API签名、覆盖率报告、编译错误日志和上一轮生成的代码时，窗口容量捉襟见肘，关键信息被挤出注意力范围的情况并不罕见\cite{liu2024lost}。此外，没有执行反馈的开环生成本质上是在"盲写"，正确性缺乏保障。本系统针对这些问题做了对应设计：多智能体架构将信息负载分散到不同角色，覆盖率反馈闭环提供执行级验证信号，迭代改进机制则让质量逐轮提升。

关于提示工程，LLM的输出质量与提示设计密切相关\cite{shah2025}，这一点在驱动生成场景中体现得尤为明显。提示需要在信息量和简洁性之间找到合适的位置：给得太少，模型缺乏必要的类型和约束信息，生成的代码编译都过不了；给得太多，注意力被稀释，模型反而忽略关键约束。本文的做法是差异化上下文注入，不同智能体根据各自职责获取不同侧面的信息。生成智能体看到的是API签名和策略描述，分析智能体看到的是覆盖率缺口和历史约束，各取所需而不互相干扰。

MAGIC4FDG采用Claude Opus作为核心模型，选型论证在第三章展开。

\section{多智能体协作框架}

多智能体系统（MAS）的核心思路是将一个复杂任务拆分给多个自主智能体协作完成。在LLM驱动的多智能体系统中，每个智能体由一个LLM实例驱动，有特定的角色定义和知识边界，彼此通过结构化消息通信。相比单模型方案，多智能体架构有三方面优势：职责分离避免单个智能体的注意力稀释；各实例可以使用不同的生成参数针对不同任务优化；模块化设计便于独立测试和替换。


协作架构的选择直接影响系统效率。常见模式包括：流水线（严格按顺序依次处理）、黑板（共享全局状态空间）、合同网（动态竞标分配任务）和DAG（按依赖关系决定执行顺序，支持条件分支和并行）。在模糊驱动生成场景中，各环节之间有明确的先后依赖：知识提取完成后才能规划策略，编译通过后才能执行模糊测试，但又不是纯线性的，中间存在条件判断和并行机会，因此DAG是最合适的架构选择。

近年来，LLM多智能体框架在软件工程领域进展很快。MetaGPT\cite{hong2023}模拟软件公司组织结构，通过产品经理、架构师、工程师等角色协作完成开发。ChatDev\cite{qian2023}模拟团队沟通模式，用对话驱动协调。微软的AutoGen\cite{wu2023autogen}采用可定制的对话式架构，在代码生成和数学推理上表现不错，但对话驱动的交互更适合探索性任务，对于有明确依赖关系的固定流程（如编译-测试-分析循环）不如DAG高效。在模糊测试领域，CKGFuzzer\cite{xu2024}引入了多智能体，但其协作只是顺序串联，各智能体共享相同上下文，缺乏差异化知识视图和基于反馈的动态调整。

LangGraph\cite{langgraph2024}是LangChain团队推出的多智能体编排框架，专门面向有状态的多参与者LLM应用。其核心抽象是StateGraph，以共享状态为中心的有向图，节点代表智能体或处理函数，边代表执行流转。本文选择LangGraph主要基于两点考虑：一是条件路由能力，可以根据当前状态动态决定下一步执行哪个节点（如编译失败则跳回修复、覆盖率未提升则提前终止）；二是状态持久化支持，任意节点都能保存和恢复执行状态，为长时间迭代任务提供断点恢复能力。

本文选择LangGraph的DAG架构作为系统的协作基础设施。模糊驱动生成的各环节（知识提取→策略规划→代码生成→编译修复→模糊测试→覆盖率分析→缺口诊断）有明确的依赖关系，天然适合DAG表达；条件路由实现自适应流程控制；状态持久化支持跨轮次知识累积和断点恢复。具体设计决策在第三章展开。

\section{静态程序分析与clang AST}

静态程序分析是在不执行程序的前提下，通过对源代码进行结构和行为建模来推断程序性质的技术\cite{Nielson1999}。其优势在于不受特定输入限制，理论上能覆盖所有可能的执行路径。常用技术包括词法分析、语法分析、类型推导、数据流分析等\cite{Bessey2010}。在本文的场景中，静态分析的价值在于从被测库源码中提取结构化的API知识，为LLM提供精确的上下文信息。

AST（抽象语法树）是静态分析中最基础的程序表示\cite{Duffy2014}。它将源代码组织为树形结构，每个节点对应一个语法单元（函数定义、变量声明、表达式等）。相比原始文本，AST消除了空白和注释等无关信息，只保留语义结构，适合做类型分析和调用关系提取。

本文的知识提取基于clang编译器前端的AST JSON导出（\texttt{clang -Xclang -ast-dump=json}）。之所以选clang而不是tree-sitter之类的轻量解析器，根本原因在于C/C++的复杂性远超表面语法。宏展开、模板实例化、typedef链解析、条件编译，这些都需要经过完整的编译语义分析才能得到正确结果，tree-sitter做不到这一点。而clang的JSON输出中，每个AST节点都带有\texttt{id}、\texttt{kind}、\texttt{type}、\texttt{loc}等字段，函数签名和参数类型可以直接读取，无需二次推断。

知识提取的具体流程如下：先从\texttt{FunctionDecl}节点中筛选出公开API签名；再通过预处理器输出识别宏常量和枚举值；然后分析\texttt{CallExpr}节点构建函数间的调用图；之后根据类型特征和命名模式对API做语义标注（区分构造/析构、读/写、配置类接口等）；最后按功能相关性对API进行分组。各阶段的实现细节在第四章展开。

关于替代方案的排除：tree-sitter的设计目标是快速增量解析以支持编辑器高亮，不经过完整的编译语义分析，无法提供宏展开后的真实结构和完整类型信息，而这些恰恰是驱动生成所必需的。libclang的Python绑定虽然提供了AST访问接口，但遍历效率低且处理复杂模板时稳定性不足。clang AST JSON导出在Docker容器内一次性完成编译和导出，既保证精确性又实现了环境隔离。

\section{本章小结}

本章介绍了MAGIC4FDG系统依赖的五项核心技术。模糊测试和模糊驱动定义了任务目标；LibFuzzer和llvm-cov提供了精确的行级覆盖率反馈；大语言模型的代码生成能力是系统的核心智能；多智能体框架提供了职责分离和协作编排的基础设施；clang AST静态分析为生成过程提供了结构化上下文知识。这五项技术共同构成了系统的技术基础，为后续的设计和实验提供支撑。